\section*{3. Kummer and Artin--Schreier theories}

\subsection*{3.0}
We denote by \((\Gm)_X\) the sheaf represented on \(X\) by the multiplicative group
\(\Spec \Oo_X[t,t^{-1}]\). Thus, for an \(\etale\) morphism \(U\to X\), its sections over \(U\) are the units
\(\Gamma(U,\Oo_U)^*\). Let \(n\in\Nn\). The kernel of the \(n\)-th power map on \((\Gm)_X\) is the sheaf of \(n\)-th roots of unity, denoted \((\mu_n)_X\). There is therefore an exact sequence at the left
\[
0\longrightarrow (\mu_n)_X\longrightarrow (\Gm)_X\xrightarrow{n}(\Gm)_X .
\tag{3.1}
\]
If \(n\) is invertible on \(X\), that is, prime to all residual characteristics of \(X\), the \(n\)-th power map is surjective as a morphism of \(\etale\) sheaves:

\begin{statement}{Kummer theory 3.2}
If \(n\) is invertible on \(X\), then
\[
0\longrightarrow (\mu_n)_X\longrightarrow (\Gm)_X\xrightarrow{n}(\Gm)_X\longrightarrow 0
\]
is exact.
\end{statement}

Indeed, for \(u\in(\Gm)_X(U)\), with \(U\to X\) \(\etale\), one must find an \(\etale\) covering on which \(u\) becomes an \(n\)-th power. Since \(n\) is invertible on \(U\), the equation
\[
T^n-u=0
\]
is separable over \(U\), so
\[
U'=\Spec\Oo_U[T]/(T^n-u)
\]
is \(\etale\) over \(U\). The morphism \(U'\to U\) is surjective, hence covering, and on \(U'\) the section \(u\) is an \(n\)-th power.

When \(n\) is invertible on \(X\), the sheaf \((\mu_n)_X\) is locally constant, in fact locally isomorphic to \((\Zz/n\Zz)_X\), and is constant in important cases. Kummer theory therefore gives information about the cohomology of \(X\) with values in certain constant sheaves. By definition
\[
H^0(X,(\Gm)_X)=\Gamma(X,\Oo_X^*),
\]
and in degree one one has the following form of Hilbert's Theorem 90.

\begin{statement}{Theorem 3.3}
There is an isomorphism
\[
H^1(X,(\Gm)_X)\xrightarrow{\sim}\Pic X,
\]
where \(\Pic X\) is the group of isomorphism classes of invertible sheaves on \(X\).
\end{statement}

Equivalently, the canonical morphism
\[
H^1(X_{\mathrm{Zar}},(\Gm)_X)\longrightarrow H^1(X_{\et},(\Gm)_X)
\]
is bijective. This amounts to saying that
\[
R^1\epsilon_*(\Gm)_X=0,
\]
where \(\epsilon:X_{\et}\to X_{\mathrm{Zar}}\) is the evident morphism of sites (VIII, 4). Thus one is reduced to the affine case. Since \(H^1\) may be computed by \v{C}ech cohomology (V, 2.5), and since for quasi-compact \(X\) it is enough to use quasi-compact coverings, the assertion follows at once from descent theory for invertible sheaves (cf. FGA, p. 190, or SGA 1, VIII, 1).

\subsubsection*{3.3.1 Cartier divisors}
Assume that \(X\) is noetherian and has no embedded components, and let \(x_j\) be the maximal points of \(X\). Let \(R_j\) be the local artinian ring of \(X\) at \(x_j\), and let \(i_j:\Spec R_j\to X\) be the inclusion. There is a canonical injection, both on \(X_{\et}\) and on \(X_{\mathrm{Zar}}\),
\[
(\Gm)_X\longrightarrow \prod_j i_{j*}(\Gm)_{\Spec R_j}.
\]
Its cokernel \(D\) is called the sheaf of Cartier divisors on \(X_{\et}\), respectively on \(X_{\mathrm{Zar}}\). In EGA IV, 21, these are simply called divisors.

\begin{statement}{Corollary 3.4}
Let \(X\) be noetherian and without embedded components, let \(\epsilon:X_{\et}\to X_{\mathrm{Zar}}\) be the canonical morphism of sites, and let \(D\) be the sheaf of Cartier divisors on \(X_{\et}\). Then \(\epsilon_*D\) is the sheaf \(D_{\mathrm{Zar}}\) of Cartier divisors on \(X_{\mathrm{Zar}}\). In particular
\[
H^0(X_{\et},D)\simeq H^0(X_{\mathrm{Zar}},D_{\mathrm{Zar}}).
\]
\end{statement}

This follows by applying \(\epsilon_*\) to the exact sequence
\[
0\longrightarrow (\Gm)_X\longrightarrow \prod_j i_{j*}(\Gm)_{\Spec R_j}\longrightarrow D\longrightarrow 0,
\]
which remains exact because \(R^1\epsilon_*(\Gm)_X=0\) by Theorem 3.3.

In characteristic \(p\ne 0\), the following result will sometimes replace Kummer theory for the study of \(p\)-torsion coefficients.

\begin{statement}{Artin--Schreier theory 3.5}
If \(X\) is of characteristic \(p>0\), there is an exact sequence
\[
0\longrightarrow (\Zz/p\Zz)_X\longrightarrow \Oo_X\xrightarrow{\wp}\Oo_X\longrightarrow 0,
\]
where \(\wp\) is the morphism of additive groups defined by \(\wp(f)=f^p-f\).
\end{statement}

The kernel of \(\wp\) is the sheaf of sections of \(\Oo_X\) which locally lie in the prime field, hence is constant. The map \(\wp\) is surjective for the \(\etale\) topology, since the equation \(T^p-T-a\), with \(a\in\Gamma(U,\Oo_U)\), is always separable in characteristic \(p\), hence defines an \(\etale\) covering of \(U\). In applying this exact sequence one uses VII, 4.2.

\begin{statement}{Proposition 3.6}
\begin{enumerate}[label=(\roman*)]
\item Let \(X\) be a scheme and let \(i:x\to X\) be the inclusion of a point. Then
\[
H^1(X,i_*(\Zz)_x)=0,
\]
where \((\Zz)_x\) denotes the constant sheaf of value \(\Zz\) on the point \(x\).
\item If \(X\) is irreducible and if, for every geometric point \(\bar x\) of \(X\), the strict localization \(X_{\bar x}\) is irreducible, in other words if \(X\) is geometrically unibranch, then
\[
H^1(X,(\Zz)_X)=0.
\]
\end{enumerate}
\end{statement}

\begin{prooftext}
For (i), the Leray spectral sequence
\[
E_2^{pq}=H^p(X,R^qi_*(\Zz)_x)\Longrightarrow H^{p+q}(x,(\Zz)_x)
\]
gives an injection \(H^1(X,i_*(\Zz)_x)\to H^1(x,(\Zz)_x)\). But, since \(x=\Spec k(x)\), the group \(H^1(x,(\Zz)_x)\) is zero: this follows from the fact that \(H^q(x,F)\), for \(q>0\), is torsion (VII, 2.2 and CG), and from the exact sequences
\[
0\longrightarrow \Zz\xrightarrow{n}\Zz\longrightarrow \Zz/n\Zz\longrightarrow 0.
\]
For (ii), let \(i:x\to X\) be the inclusion of the generic point. Under the stated hypotheses the canonical morphism \((\Zz)_X\to i_*(\Zz)_x\) is immediately seen to be bijective, and the assertion follows from (i). The converse to this bijectivity is also true.
\end{prooftext}

\begin{remarktext}{Remarks 3.7}
In Proposition 3.6(i) and (ii), \(\Zz\) may be replaced by any torsion-free abelian group. On the other hand, (ii) becomes false if the geometrically unibranch hypothesis is dropped; for example take an algebraic curve over a field with an ordinary double point and compute \(H^1(X,\Zz)\).
\end{remarktext}

\section*{4. The case of an algebraic curve}

\subsection*{4.0}
Let \(X\) be a noetherian scheme of dimension one. Let \(R(X)\) be the ring of rational functions on \(X\), and let \(i:\Spec R(X)\to X\) be the inclusion. The ring \(R(X)\) is artinian: it is the product of the local rings of \(X\) at its maximal points, so one may use the results of VIII, 2.

Let \(F\) be an abelian sheaf on \(\Spec R(X)\), and consider the sheaves \(R^qi_*F\), \(q>0\), on \(X\). It follows at once from VIII, 5.2 that the fibre of \(R^qi_*F\) at a geometric point lying over a maximal point of \(X\) is zero. Since \(X\) has dimension one, such a sheaf is of a very special kind.

\begin{statement}{Lemma 4.1}
Let \(X\) be noetherian and let \(F\) be a sheaf on \(X\). The following conditions are equivalent; a sheaf satisfying them will be called a skyscraper sheaf.
\begin{enumerate}[label=(\roman*)]
\item The fibre of \(F\) at every geometric point \(\bar x\) lying over a non-closed point \(x\in X\) is zero.
\item Each section \(s\in F(X')\), with \(X'\to X\) \(\etale\) of finite type, is zero except at finitely many closed points of \(X'\).
\item One has an isomorphism
\[
F\simeq \bigoplus_x i_{x*}i_x^*F,
\]
where \(x\) runs through the closed points of \(X\), and \(i_x:x\to X\) is the inclusion.
\end{enumerate}
\end{statement}

\begin{prooftext}
If (i) holds and \(z\in F(X')\), its support is a closed subset of \(X'\) formed of closed points, hence is finite because \(X'\) is noetherian; this proves (ii). Conversely, (ii) implies (i) by the calculation of fibres (VIII, 3.9). Condition (iii) plainly implies (i), since the functor of fibres commutes with direct sums and the support of \(i_{x*}G_x\) is contained in \(\{x\}\). Finally, assuming (ii), the canonical morphism
\[
F\longrightarrow \prod_x i_{x*}i_x^*F
\]
factors through the direct sum; the resulting morphism to \(\bigoplus_x i_{x*}i_x^*F\) is checked directly to be bijective.
\end{prooftext}

\subsubsection*{4.1.1}
Applying the lemma to the sheaf \(R^qi_*F\) introduced above gives
\[
R^qi_*F\simeq \bigoplus_x i_{x*}i_x^*(R^qi_*F),
\]
with \(x\) running through the closed points of \(X\). By VIII, 5, the fibre of \(R^qi_*F\) at a geometric point \(\bar x\) of \(X\) is
\[
(R^qi_*F)_{\bar x}\simeq H^q(\Spec R(X_{\bar x}),F),
\]
where \(X_{\bar x}\) is the strict localization of \(X\) at \(\bar x\), and \(F\) also denotes the induced sheaf on \(\Spec R(X_{\bar x})\). Assume now that the residue field of every closed point of \(X\) is separably closed. Then
\[
H^p(X,R^qi_*F)\simeq
\begin{cases}
\displaystyle\bigoplus_x H^q(\Spec R(X_x),F),&p=0,\\
0,&p>0,
\end{cases}
\]
for \(q>0\), where \(X_x\) denotes the strict localization at \(x\).

\begin{statement}{Corollary 4.2}
Let \(X\) be a noetherian scheme of dimension one such that, for every closed point \(x\), the residue field \(k(x)\) is separably closed. Let \(F\) be an abelian sheaf on \(\Spec R(X)\). In the Leray spectral sequence
\[
E_2^{pq}=H^p(X,R^qi_*F)\Longrightarrow H^{p+q}(\Spec R(X),F)
\]
one has
\[
E_2^{pq}=0\quad(p>0, q>0),\qquad
E_2^{0q}=\bigoplus_x H^q(\Spec R(X_x),F)\quad(q>0).
\]
Consequently there is an exact sequence
\[
\begin{aligned}
0&\to H^1(X,i_*F)\to H^1(\Spec R(X),F)\to\bigoplus_x H^1(\Spec R(X_x),F)\to H^2(X,i_*F)\to{}\\
&\to H^2(\Spec R(X),F)\to\bigoplus_x H^2(\Spec R(X_x),F)\to\cdots .
\end{aligned}
\]
\end{statement}

\subsection*{4.3}
Assume further that the \(\ell\)-cohomological dimension of \(\Spec R(X)\) is at most one. Let \(K^1,\ldots,K^m\) be the residue fields of the artinian ring \(R(X)\), equivalently the residue fields at the maximal points of \(X\). By the dictionary of VIII, 2,
\[
\operatorname{cd}_\ell(\Spec R(X))=\sup_i\operatorname{cd}_\ell(G(\overline K^i/K^i)).
\]
Every residue field \(K_x\) of \(R(X_x)\) is a separable extension, possibly infinite, of one of the \(K^i\), hence
\[
\operatorname{cd}_\ell(G(\overline K_x/K_x))\leq \operatorname{cd}_\ell(G(\overline K^i/K^i))
\]
by CG II, 4.1, Proposition 10. Applying Corollary 4.2 to an \(\ell\)-torsion sheaf \(F\), one obtains
\[
\begin{cases}
0\to H^1(X,i_*F)\to H^1(\Spec R(X),F)\to \displaystyle\bigoplus_x H^1(\Spec R(X_x),F)\to H^2(X,i_*F)\to 0,\\
H^q(X,i_*F)=0\quad(q>2).
\end{cases}
\tag{4.3.1}
\]
The hypothesis \(\operatorname{cd}_\ell(\Spec R(X))\leq 1\) is satisfied, for example, when \(X\) is an algebraic curve over a separably closed field, by Tsen's theorem. Exact sequences of this type were studied by Ogg and Shafarevich.

\subsection*{4.4}
Assume that \(\ell\) is invertible on \(X\), and take \(F=(\Gm)_{R(X)}\). From Theorem 3.3 and CG, Chapter I, Proposition 12, one has, under the same hypothesis \(\operatorname{cd}_\ell(\Spec R(X))\leq 1\),
\[
H^1(\Spec R(X),\Gm)=0,
\qquad H^q(\Spec R(X),\Gm)\text{ has no }\ell\text{-torsion for }q\geq2.
\tag{4.4.1}
\]
The same holds with \(R(X)\) replaced by \(R(X_x)\). Corollary 4.2 therefore gives
\[
H^1(X,i_*(\Gm)_{R(X)})=0,
\qquad H^q(X,i_*(\Gm)_{R(X)})\text{ has no }\ell\text{-torsion for }q>1.
\tag{4.4.2}
\]
For any sheaf \(F\) on \(X\), Lemma 4.1 shows that the kernel and cokernel of the canonical homomorphism \(F\to i_*i^*F\) are skyscraper sheaves. Since the closed residue fields are separably closed, these skyscraper sheaves have no higher cohomology. Thus
\[
H^i(X,F)\longrightarrow H^i(X,i_*i^*F)
\]
is an isomorphism for \(i\geq2\). Applied to \((\Gm)_X\), this yields
\[
H^i(X,\Gm)\text{ has no }\ell\text{-torsion for }i\geq2,
\tag{4.5}
\]
under the hypotheses of (4.4). Combining this with Kummer theory and Theorem 3.3 gives:

\begin{statement}{Theorem 4.6}
Let \(X\) be noetherian of dimension one, let \(n\) be an integer invertible on \(X\), suppose that \(\operatorname{cd}_\ell(R(X))\leq1\) for every prime \(\ell\) dividing \(n\), and suppose that for every closed point \(x\) of \(X\) the residue field \(k(x)\) is separably closed. Then
\[
H^q(X,(\mu_n)_X)=0\quad(q>2),
\]
and for \(q\leq2\) the cohomology is given by the exact sequence
\[
\begin{aligned}
0\to H^0(X,(\mu_n)_X)&\to \Gamma(X,\Oo_X^*)\xrightarrow{n}\Gamma(X,\Oo_X^*)\to H^1(X,(\mu_n)_X)\\
&\to \Pic X\xrightarrow{n}\Pic X\to H^2(X,(\mu_n)_X)\to0.
\end{aligned}
\]
\end{statement}

\begin{statement}{Corollary 4.7}
Let \(X\) be a complete connected curve over a separably closed field \(k\), of characteristic prime to \(n\). Then
\[
H^0(X,(\mu_n)_X)\simeq\mu_n(k),
\]
\[
H^1(X,(\mu_n)_X)\simeq {}_n\Pic X,
\]
and
\[
H^2(X,(\mu_n)_X)\simeq (\Pic X)/n\simeq (\Zz/n\Zz)^c,
\]
where \(c\) is the number of irreducible components of \(X\).
\end{statement}

For the last assertion, recall that the Picard scheme of such an \(X/k\) has a decomposition
\[
0\longrightarrow A\longrightarrow \mathcal{P}ic_{X/k}\longrightarrow \Zz^c\longrightarrow 0,
\tag{4.8}
\]
where \(A\) is a connected algebraic group over \(k\). Multiplication by \(n\) on such an \(A\) is \(\etale\), hence surjective on \(k\)-valued points. Thus \({}_n(\Pic X)\simeq {}_nA(k)\) and \((\Pic X)/n\simeq(\Zz/n\Zz)^c\). If \(X\) is smooth over \(k\), then \(A\) is the Jacobian and \({}_n\Pic X\simeq(\Zz/n\Zz)^{2g}\), where \(g\) is the genus. Since \(\mu_n(k)\simeq\Zz/n\Zz\), the ranks of the cohomology groups are \(1,2g,1\), as expected.

\section*{5. The trace method}

\subsection*{5.1}
Let \(X\) be a scheme, \(f:X'\to X\) a finite \(\etale\) covering, and \(F\) a sheaf on \(X\). Adjunction gives a restriction morphism
\[
F\xrightarrow{\mathrm{res}} f_*f^*F.
\]
Since \(f\) is finite, \(H^q(X,f_*f^*F)\simeq H^q(X',f^*F)\), whence the usual restriction map
\[
H^q(X,F)\xrightarrow{\mathrm{res}} H^q(X',f^*F).
\tag{5.1.1}
\]
There is also a trace morphism
\[
f_*f^*F\xrightarrow{\mathrm{tr}}F,
\tag{5.1.2}
\]
which induces
\[
H^q(X',f^*F)\xrightarrow{\mathrm{tr}}H^q(X,F).
\tag{5.1.3}
\]
It is characterized by being local for the \(\etale\) topology and, when \(X'\) is the disjoint sum of \(d\) copies of \(X\), by being the summation map \(F^d\to F\). The uniqueness is immediate, since every finite \(\etale\) covering is locally constant; existence follows by the usual descent argument. Hence
\[
\mathrm{tr}\circ\mathrm{res}:F\to F
\]
is multiplication by the local degree of \(f\). If the degree is a constant \(d\), the same is true on cohomology.

\begin{statement}{Corollary 5.2}
Let \(f:X'\to X\) be a finite \(\etale\) covering of constant degree \(d\). If \(F\) is a sheaf of \(r\)-torsion on \(X\), with \((r,d)=1\), then \(H^q(X',f^*F)=0\) implies \(H^q(X,F)=0\).
\end{statement}

Indeed, multiplication by \(d\) is an isomorphism of \(F\), hence of its cohomology, and the resulting isomorphism factors through zero.

\subsubsection*{5.2.1}
Let \(i:U\to X\) be an immersion, let \(f:U'\to U\) be a finite \(\etale\) covering of degree \(d\), and suppose that \(f\) is induced by a finite morphism \(g:X'\to X\) by base change:
\[
\xymatrix@C=2cm@R=1.25cm{
U'\ar[r]^{i'}\ar[d]_f&X'\ar[d]^g\\
U\ar[r]^i&X.
}
\]
For a sheaf \(F\) on \(U\) one has the composites
\[
F\xrightarrow{\mathrm{res}}f_*f^*F\xrightarrow{\mathrm{tr}}F,
\]
whose composite is multiplication by \(d\). Since \(i_*f_*=g_*i'_*\), this gives
\[
i_*F\xrightarrow{\mathrm{res}}g_*i'_*f^*F\xrightarrow{\mathrm{tr}}i_*F,
\tag{5.3}
\]
and since \(i_!f_*\simeq g_*i'_!\),
\[
i_!F\xrightarrow{\mathrm{res}}g_*i'_!f^*F\xrightarrow{\mathrm{tr}}i_!F,
\tag{5.4}
\]
again with composite multiplication by \(d\). Passing to cohomology gives the corresponding restriction--trace composites for \(i_*F\) and \(i_!F\).

\begin{statement}{Proposition 5.5}
Let \(X\) be quasi-compact and quasi-separated, let \(\ell\) be a prime, and let \(H\) be a semi-exact functor to \(\mathbf{Ab}\), defined on the category of \(\ell\)-sheaves on \(X\), and compatible with filtered inductive limits. The following are equivalent:
\begin{enumerate}[label=(\roman*)]
\item \(H=0\).
\item \(H(f_*i_!(\Zz/\ell\Zz))=0\) for every finite morphism \(f:Y\to X\), and every open immersion \(i:U\to Y\) with \(U\) of finite presentation over \(X\). If \(X\) is noetherian, one may restrict to integral \(Y\).
\end{enumerate}
\end{statement}

\begin{prooftext}
By Corollary 2.7.2 and Proposition 2.5 it is enough to test \(H\) on sheaves \(i_!G\), where \(i:U\to X\) is the inclusion of a subscheme of finite presentation over \(X\), and \(G\) is an isotrivial locally constant \(\ell\)-sheaf on \(U\). Choose a principal covering \(U''\to U\), with group \(\mathfrak g\), on which \(G\) becomes constant. Let \(A\) be the \(\ell\)-group of its values, and let \(\mathfrak h\) be an \(\ell\)-Sylow subgroup of \(\mathfrak g\). If \(U'=U''/\mathfrak h\), then the degree \(d\) of \(U'\to U\) is prime to \(\ell\). Extending \(U'\to U\) to a finite morphism \(X'\to X\) (EGA IV, 8.12.6), and applying (5.4), one reduces to proving the vanishing for the pullback of \(G\) to \(U'\). As an \(\mathfrak h\)-module, \(A\) admits a filtration whose successive quotients are \(\Zz/\ell\Zz\) with trivial action; the assertion follows by semi-exactness.
\end{prooftext}

\begin{statement}{Proposition 5.6}
Let \(X\) be noetherian, let \(\ell\) be a prime, and let \(\varphi\) be a function with values in \(\Nn\), defined on the integral schemes finite over \(X\). Suppose \(\varphi(Y_1)\leq\varphi(Y_2)\) whenever there is an \(X\)-morphism \(Y_1\to Y_2\), and that the inequality is strict if \(Y_1\) is a proper closed subscheme of \(Y_2\). For finite \(Y\) over \(X\), put \(\varphi(Y)=\sup\varphi(Y_i)\), where the \(Y_i\) are the irreducible components of \(Y\), with their reduced induced structures. The following are equivalent:
\begin{enumerate}[label=(\roman*)]
\item For every finite \(Y\) over \(X\) and every \(\ell\)-sheaf \(F\) on \(Y\), one has \(H^q(Y,F)=0\) for \(q>\varphi(\operatorname{Supp}F)\).
\item For every integral finite \(Y\) over \(X\), one has \(H^q(Y,\Zz/\ell\Zz)=0\) for \(q>\varphi(Y)\).
\end{enumerate}
\end{statement}

\begin{prooftext}
Assuming (ii), apply Proposition 5.5 to each \(H^q\), using the facts that cohomology commutes with inductive limits and with direct image under finite morphisms. One is reduced to the case where \(Y\) is integral finite over \(X\), and \(F=i_!(\Zz/\ell\Zz)_U\) for a non-empty open immersion \(i:U\to Y\). The exact sequence
\[
0\longrightarrow i_!(\Zz/\ell\Zz)_U\longrightarrow (\Zz/\ell\Zz)_Y\longrightarrow (\Zz/\ell\Zz)_Z\longrightarrow0,
\]
with \(Z=Y-U\ne Y\), and induction on \(\varphi\), give the vanishing. The converse is immediate.
\end{prooftext}

\begin{statement}{Corollary 5.7}
Let \(X\) be an algebraic curve over a separably closed field, of characteristic different from \(\ell\). Then the cohomology of \(X\) with values in any \(\ell\)-torsion sheaf \(F\) vanishes in degrees \(>2\). If \(X\) is affine, it vanishes in degrees \(>1\).
\end{statement}

For finite integral \(Y\to X\), put \(\varphi(Y)=2\dim Y\), or \(\varphi(Y)=\dim Y\) in the affine case. Proposition 5.6 reduces the assertion to the case \(F=\Zz/\ell\Zz\) on an integral curve. Since \((\Zz/\ell\Zz)_Y\simeq(\mu_\ell)_Y\), Theorem 4.6 gives the vanishing in degrees \(>2\). In the affine case one may pass to the normalization and then to a smooth compactification \(\bar X\); at least one point is missing from \(\bar X\), so \(A(k)\to\Pic X\) is surjective, where \(A\) is the Jacobian of \(\bar X\). As \(A(k)\) is divisible by \(\ell\ne\operatorname{char}k\), so is \(\Pic X\); hence \(H^2(X,\mu_\ell)=0\) by Theorem 4.6.

\begin{remarktext}{Remark 5.8}
For later reference, the proof of Proposition 5.5 actually proves the following. Let \(X\) be quasi-compact and quasi-separated, \(\ell\) a prime, and \(C\) a strictly full subcategory of the category of constructible abelian \(\ell\)-torsion sheaves, stable under direct summands and extensions. Suppose that for every finite morphism \(f:Y\to X\) and every open immersion \(i:U\to Y\) with \(U\) of finite presentation over \(X\), one has
\[
f_*(i_!(\Zz/\ell\Zz)_U)\in C.
\]
Then \(C\) contains all constructible \(\ell\)-torsion sheaves.
\end{remarktext}

\subsection*{Bibliography}
\begin{enumerate}[label={[\arabic*]}]
\item P. Gabriel, \emph{Des Cat\'egories Ab\'eliennes}, thesis, Paris, 1962.
\item A. Ogg, Cohomology of abelian varieties over function fields, \emph{Annals of Mathematics} 76, no. 2 (1962).
\item M. Deuring, Die Typen der Multiplikatorenringe elliptischer Funktionenk\"orper, \emph{Hamb. Abh.} 13 (1940), p. 197.
\item A. Grothendieck, \emph{Fondements de la g\'eom\'etrie alg\'ebrique}, excerpts from S\'eminaire Bourbaki 1957--1962, Secr\'etariat math\'ematique, Paris. Cited as FGA.
\item J.-P. Serre, \emph{Cohomologie galoisienne}, Lecture Notes in Mathematics 5, Springer. Cited as CG.
\item A. Douady, Cohomologie des groupes compacts totalement discontinus, S\'eminaire Bourbaki, expos\'e 189.
\end{enumerate}
